\documentclass[11pt,a4paper]{article}
\usepackage{fontspec}
\setmainfont{PT Serif}[Ligatures=TeX]
\setsansfont{Helvetica Neue}[Ligatures=TeX]
\setmonofont{Menlo}
\usepackage{polyglossia}
\setdefaultlanguage{russian}
\setotherlanguage{english}
\newfontfamily\cyrillicfont{PT Serif}[Ligatures=TeX]
\newfontfamily\cyrillicfontsf{Helvetica Neue}
\newfontfamily\cyrillicfonttt{Menlo}
\usepackage{amsmath,amssymb,amsthm}
\usepackage{geometry}
\geometry{margin=22mm}
\usepackage{booktabs}
\usepackage{makecell}
\renewcommand{\cellalign}{tc}\renewcommand{\theadalign}{tc}\renewcommand{\theadfont}{\small}
\usepackage{graphicx}
\usepackage{enumitem}
\setlist{nosep,leftmargin=*}
\usepackage{microtype}
\usepackage{hyperref}
\hypersetup{colorlinks,linkcolor=black,urlcolor=blue!50!black}
\renewcommand{\arraystretch}{1.12}
\emergencystretch=2em
\newcommand{\E}{\mathbb{E}}
\newcommand{\Prob}{\mathbb{P}}
\newcommand{\Bern}{\operatorname{Bern}}
\newcommand{\Bin}{\operatorname{Bin}}
\newcommand{\tb}[1]{\resizebox{\linewidth}{!}{#1}}
\theoremstyle{plain}
\newtheorem{proposition}{Утверждение}
\newtheorem{theorem}{Теорема}
\theoremstyle{definition}
\newtheorem{definition}{Определение}
\newtheorem{remark}{Замечание}
\newtheorem{lemma}{Лемма}

\title{Задача B «Экономия на эвалах»: решение\\[2pt]\large ограничение на долю пропущенных релизов с проблемами}
\author{для организаторов}
\date{25 октября 2026}

\begin{document}
\sloppy
\maketitle

\section*{Что решаем}

На релизе выполняются все запланированные прогоны: тест $i$ --- $n_i$ раз. Если хоть один прогон красный, релиз уходит на доработку; если все зелёные, релиз выпускается. Плохо, если выпущен релиз с~проблемой --- в~нём есть проблема хотя бы в~одном тесте. Требование организаторов --- на вероятность пропуска \emph{при условии, что релиз проблемный}:
\[
\pi_{\rm c}(n)=\Prob(\text{все прогоны плана зелёные}\mid\text{релиз с проблемой})=\frac{\pi(n)}{P_1}\le\alpha=0.1,
\]
где $\pi(n)=\Prob(\text{релиз с проблемой, а все прогоны плана зелёные})$ и $P_1=1-\prod_i(1-p_i)$, с вероятностью не меньше $0.9$ по истории; стоимость плана --- $\sum_ic_in_i$. В мире задачи $P_1=0.81$, так что в безусловных терминах требование --- $\pi(n)\le\alpha P_1=0.081$. Команды $P_1$ не знают и оценивают сами: по модели ($1-\prod(1-\hat p_i)=0.84$ на опубликованной истории) или по доле релизов истории с красным ($56/70=0.80$). Командам устно названо «5--10\%»; на валидации считается число пропущенных релизов среди проблемных и стоимость.

\paragraph{Схема решения.}
\begin{enumerate}
\item Один релиз: формула для $\pi(n)$ (раздел~1).
\item Известные параметры: после логарифма --- рюкзак; пороговое правило (раздел~2).
\item Что можно оценить по данным: видимая часть $\pi$ оценивается несмещённо, невидимая --- никак (раздел~3).
\item Оценки по тесту: ОММ, ОМП, частоты; в этом мире событий мало, и это главная трудность (раздел~4).
\item Погрешность при фиксированном плане (раздел~5).
\item Контроль ожидаемого числа упущенных проблем вместо $\pi$ (раздел~6).
\item Отбор плана по тем же данным: бутстреп всей процедуры и точная проверка на отложенных релизах (раздел~7).
\item Мир задачи, результаты, оценивание ответов (разделы~8--10).
\end{enumerate}

\section*{Обозначения}

\begin{center}\small
\begin{tabular}{@{}p{0.36\linewidth}p{0.6\linewidth}@{}}
\toprule
$i=1,\dots,M$, $M=40$ & тесты \\
$c_i$ & стоимость одного прогона теста $i$ \\
$p_i$ & вероятность того, что в сценарии теста $i$ на релизе есть проблема \\
$s_i$ & вероятность того, что прогон покраснеет, если проблема есть (без проблемы прогон всегда зелёный) \\
$n_i\in\{0,\dots,5\}$, $n$ & сколько раз гонять тест $i$; вектор $n$ --- \emph{план} \\
$g_i(m)=1-p_i+p_i(1-s_i)^m$ & вероятность того, что первые $m$ прогонов теста $i$ зелёные \\
$\pi(n)$, $\pi_{\rm c}(n)=\pi(n)/P_1$ & вероятность пропустить релиз с проблемой при плане $n$: по всем релизам и при условии, что релиз проблемный \\
$P_0=\prod_i(1-p_i)$, $P_1=1-P_0$ & вероятности чистого и проблемного релиза \\
$K_{ij}\in\{0,\dots,5\}$ & число красных среди пяти записей теста $i$ на релизе $j$ истории \\
$C(n)=\sum_ic_in_i$ & стоимость плана \\
$\mathbf5$ & план «каждый тест $5$ раз» \\
$\alpha=0.1$, $\gamma=0.1$ & допустимая доля пропущенных среди проблемных релизов; допустимая вероятность того, что правило ошибётся \\
$\bar\varepsilon=0.02$ & предположение: $\pi(\mathbf5)\le\bar\varepsilon$ \\
$N_{\rm h}=70$ & релизов в истории; на каждом все тесты прогнаны $5$ раз \\
\bottomrule
\end{tabular}
\end{center}

\section{Один релиз}

Тесты независимы, поэтому все прогоны плана зелёные с вероятностью $\prod_ig_i(n_i)$. Из этого события вычитаем чистый релиз (проблем нет --- все прогоны зелёные), его вероятность $\prod_i(1-p_i)$:
\[
\boxed{\ \pi(n)=\prod_ig_i(n_i)-\prod_i(1-p_i).\ }
\]
В отличие от ожидаемого числа упущенных проблем, $\pi$ не складывается по тестам: релиз с двумя проблемами пропускается, только если упущены обе. Поэтому ловить нужно не каждую проблему, а хотя бы одну на релизе.

\section{Если бы параметры были известны}

\[
\min_n\ \sum_ic_in_i\qquad\text{при}\qquad\pi(n)\le\alpha P_1 .
\]
Ограничение равносильно $\prod_ig_i(n_i)\le\alpha P_1+P_0=\alpha+(1-\alpha)P_0$. После логарифма
\[
\sum_iD_i(n_i)\ \ge\ B,\qquad D_i(n)=-\log g_i(n),\qquad B=-\log\bigl(\alpha+(1-\alpha)P_0\bigr)
\]
--- слева сумма по тестам, справа константа. Это рюкзак: $k$-й прогон теста $i$ --- предмет со стоимостью $c_i$ и пользой
\[
d_i(k)=D_i(k)-D_i(k-1)=-\log\bigl(1-r_i(k)\bigr),\qquad r_i(k)=\frac{p_is_i(1-s_i)^{k-1}}{1-p_i+p_i(1-s_i)^{k-1}},
\]
где $r_i(k)$ --- вероятность того, что $k$-й прогон покраснеет, если предыдущие прогоны теста зелёные.

\begin{lemma}
$d_i(1)\ge d_i(2)\ge\dots\ge d_i(5)$.
\end{lemma}

\begin{proof}
При $x=(1-s_i)^{k-1}$ величина $r_i(k)=p_is_ix/(1-p_i+p_ix)$ возрастает по $x$, а $x$ убывает по $k$.
\end{proof}

\begin{theorem}
Если брать прогоны по убыванию $d_i(k)/c_i$, пока не выполнится $\sum_iD_i(n_i)\ge B$, получится план не дороже оптимального плюс стоимость одного прогона: $C(n^{\rm gr})\le C^*+\max_ic_i$.
\end{theorem}

\begin{proof}
Непрерывная релаксация --- рюкзак с $x_{ik}\in[0,1]$; её оптимум --- брать предметы по убыванию $d/c$, последний частично. По лемме внутри теста $d_i(k)/c_i$ убывает, поэтому $k$-й прогон берётся только после $(k-1)$-го, и решение релаксации --- план стоимостью не больше $C^*$. Жадный план отличается последним предметом, взятым целиком.
\end{proof}

\paragraph{Пороговое правило.} При пороге $\lambda$ тест гоняется $k$-й раз, если $d_i(k)/c_i\ge\lambda$: вероятность того, что прогон поймает проблему, в расчёте на рубль выше порога. Почти неломающиеся тесты не гоняются, точные --- один раз, часто ломающиеся с умеренной $s_i$ --- несколько раз. Точный оптимум --- целочисленное линейное программирование с ограничениями $\sum d_i(k)x_{ik}\ge B$, $x_{i,k}\le x_{i,k-1}$. Жадный план в мирах задачи дороже точного в среднем на $1.1\%$, в худшем из $20$ миров --- на $6\%$ (оценка $C^*+\max_ic_i$ грубая: самый дорогой прогон --- пятая часть стоимости оракула). Оракул стоит в среднем $0.85$ от нынешней практики «все по разу» (у неё $\pi_{\rm c}\approx0.25$) и $0.17$ от «все по $5$»; в мире $24$ --- $0.91$ и $0.18$.

\section{Что можно оценить по данным}

На релизе истории записаны все пять прогонов каждого теста, поэтому для плана $n$ видно событие
\[
V_j(n)=\{\text{первые $n_i$ записей каждого теста зелёные, но среди всех записей есть красный}\}
\]
--- «план пропустил бы релиз, а пятикратный прогон поймал бы проблему».

\begin{lemma}[разложение]
$\rho(n)=\Prob(V(n))=\prod_ig_i(n_i)-\prod_ig_i(5)$, и $\pi(n)=\rho(n)+\pi(\mathbf5)$.
\end{lemma}

\begin{proof}
Событие «зелёные все пять записей всех тестов» вложено в «зелёные все прогоны плана», разность --- $V(n)$. Кроме того, $\pi(\mathbf5)=\prod_ig_i(5)-\prod_i(1-p_i)$.
\end{proof}

\paragraph{Несмещённая оценка видимой части.} Пусть $h(k,m)=\binom{5-m}{k}\big/\binom5k$ --- вероятность того, что $m$ случайных записей из пяти зелёные, если красных $k$ (при $k=0$ это $1$). Тогда
\[
\hat\rho(n)=\frac1{N_{\rm h}}\sum_j\Bigl[\prod_ih(K_{ij},n_i)-\mathbb 1\{K_{\cdot j}=0\}\Bigr]
\]
несмещённо оценивает $\rho(n)$: тесты независимы, а $\E\,h(K,m)=\Prob(\text{первые $m$ прогонов зелёные})$.

\begin{proposition}[невидимую часть оценить нельзя]
Ни $p_i$, ни $s_i$, ни $\pi(n)$ не имеют несмещённых оценок ни при каком числе релизов.
\end{proposition}

\begin{proof}
По одному тесту достаточная статистика --- мультиномиальный вектор частот $K=0,\dots,5$ с вероятностями $P_0,\dots,P_5$; несмещённо оцениваются только многочлены от них. При $k\ge1$ в $P_k$ есть множитель $s$, а $P_0=1-\sum_{k\ge1}P_k$, поэтому любой такой многочлен при $s\to0$ стремится к константе, а $1-g(m)=p\bigl(1-(1-s)^m\bigr)\to0$ при любом $p$. Значит, $\pi$ при $s_1\to0$ зависит от $p_1$, а данные --- нет. Для $s$ --- симметрично при $p\to0$: данные вырождаются в $K\equiv0$ при любом $s$.
\end{proof}

Поэтому «пяти прогонов хватает» входит в постановку предположением $\pi(\mathbf5)\le\bar\varepsilon=0.02$ (в мирах генератора в среднем $0.012$, максимум $0.014$ на $60$ мирах), и требование сводится к видимой части: $\rho(n)\le\alpha P_1-\bar\varepsilon$ ($=0.061$ в мире задачи).

\section{Оценки по тесту}

Для теста на релизе $\Prob(K=0)=1-p+p(1-s)^5$, $\Prob(K=k)=p\binom5ks^k(1-s)^{5-k}$ при $k\ge1$.
\begin{itemize}
\item \textbf{ОММ.} Средние по релизам величин $K/5$ и $K(K-1)/20$ --- $a$ и $b$ --- несмещённо оценивают $ps$ и $ps^2$; $\hat s=b/a$, $\hat p=a^2/b$. Если у теста нет релиза с двумя красными ($b=0$), берётся общая по тестам $\hat s=\sum b/\sum a$.
\item \textbf{ОМП (EM).} На релизе с $K=0$ апостериорная вероятность проблемы $q=pz/(1-p+pz)$, $z=(1-s)^5$; шаг: $W=n_{K\ge1}+n_{K=0}\,q$, $p\leftarrow W/N_{\rm h}$, $s\leftarrow\sum K/(5W)$.
\item \textbf{Частоты.} Без модели --- оценка $\hat\rho$ выше.
\end{itemize}

\paragraph{Главная трудность.} В мире задачи $80\%$ релизов проблемные, но тестов $40$, поэтому у отдельного теста $p_i\approx0.02$--$0.06$. В опубликованной истории у типичного теста $2$ релиза с красным, у $9$ тестов --- ни одного, у $21$ --- не больше двух. Отсюда:
\begin{itemize}
\item ОМП часто уходит на границу: если все релизы с проблемой дали по одному красному, $\hat p\to1$, $\hat s\to0$, и тест выглядит «ломается всегда, но не видно»; такой план гоняет всё;
\item ОММ делит на малую и шумную $b$ и в среднем завышает $\pi$ --- это перестраховка;
\item частоты несмещённые, но у теста без событий вклад ровно $0$ --- как и у ОММ и ОМП ($\hat p=0$ при $a=0$): ни одна оценка по данным не гоняет пять реально ломающихся тестов без событий (в мире $24$ это тесты $0,3,4,12,30$; пять таких тестов из девяти --- типично, ожидается $3.3$).
\end{itemize}

Последнее --- потолок для всех методов по данным, а не ловушка для наивных: оракул гоняет три из этих тестов, и их обнуление поднимает $\pi$ с $0.100$ до $0.126$ при экономии $0.009$ от полной проверки. Лечится общим априорным распределением по тестам (пулинг, эмпирический Байес) или одним прогоном дешёвых тестов без событий.

\section{Погрешность при фиксированном плане}

Погрешность оценки $\pi$ для заданного плана считается дельта-методом (для ОММ --- по выборочной ковариации пар $(K/5,K(K-1)/20)$ каждого теста; для ОМП --- по информации Фишера одного релиза), бутстрепом по релизам или, для частот, как погрешность среднего по релизам.

\begin{center}\small
\begin{tabular}{@{}llcccc@{}}
\toprule
план & оценка & смещение & разброс & \thead{оценённая погрешность\\/ настоящая} & \thead{покрытие\\границы $90\%$} \\
\midrule
все дважды & ОММ & $+26\%$ & $20\%$ & $0.92$ & $0.99$ \\
 & ОМП & $+168\%$ & $86\%$ & --- & $1.00$ \\
 & частоты${}^*$ & $0\%$ & $25\%$ & $0.96$ & $0.84$ \\
план оракула & ОММ & $+20\%$ & $19\%$ & $0.97$ & $0.98$ \\
 & ОМП & $+146\%$ & $75\%$ & --- & $0.99$ \\
 & частоты${}^*$ & $0\%$ & $29\%$ & $1.01$ & $0.84$ \\
\bottomrule
\end{tabular}
\end{center}

${}^*$\,оценивают видимую часть $\rho$. $10$ миров по $60$ историй. ОМП вырождается на тестах с малым числом событий, и его погрешность по информации Фишера бессмысленна. ОММ перестраховывается; частоты несмещены, но при малом числе событий сумма скошена, и нормальная граница недопокрывает.

\section{Контроль ожидаемого числа упущенных проблем вместо $\pi$}

Естественная подмена --- требовать $\E N(n)=\sum_ip_i(1-s_i)^{n_i}\le\alpha P_1$ или $\Prob(N\ge1)\le\alpha P_1$ (или те же величины, делённые на $P_1$, не больше $\alpha$), где $N$ --- число упущенных проблем. Так как пропуск релиза означает, что упущены все его проблемы,
\[
\pi(n)\le\Prob(N(n)\ge1)\le\E N(n),
\]
и обе подмены консервативны. В этом мире они очень дороги: релиз часто содержит несколько проблем, и поймать одну из них достаточно, а $\E N$ требует ловить каждую. В мире задачи $\E N(\mathbf5)\approx0.06$ сравнимо с $\alpha P_1=0.08$, поэтому план с $\E N\le\alpha P_1$ совпадает с полной проверкой.

\section{Отбор плана по тем же данным}

\paragraph{Проклятие победителя.} План выбирается по тем же оценкам, и среди дешёвых планов чаще проходит тот, у которого оценка случайно занижена. Запас, посчитанный для фиксированного плана (дельта-метод), этого не учитывает.

\paragraph{Бутстреп всей процедуры.}
\begin{enumerate}
\item перевыбрать $70$ релизов истории с возвращением;
\item на перевыборке заново оценить параметры и построить жадный путь планов;
\item для каждого запаса $c$ из сетки взять самый дешёвый план с оценкой $\pi$ не больше $\alpha-c$ и проверить, укладывается ли он в $\alpha$ по оценкам на \emph{полной} истории;
\item повторить $B=40$ раз и взять наименьший $c$, при котором план укладывается в $90\%$ повторов;
\item на полной истории взять самый дешёвый план с оценкой не больше $\alpha-c$.
\end{enumerate}
Запас учитывает и шум оценок, и отбор; гарантия приближённая.

\paragraph{Проверка на отложенных релизах (точно).}
\begin{enumerate}
\item История случайно делится на обучающую ($35$ релизов) и проверочную ($m=35$) части.
\item По обучающей --- оценки и жадный путь планов $\mathbf0=n^{(0)}\le\dots\le n^{(K)}=\mathbf5$.
\item По проверочной: $S_k$ --- число релизов, которые план $n^{(k)}$ пропустил бы (событие $V$). Идём от $\mathbf5$ к дешёвым планам, пока верхняя граница Клоппера--Пирсона $U_\gamma(S_k,m)\le\alpha\hat P_1-\bar\varepsilon$, где $\hat P_1$ --- оценка доли проблемных релизов по обучающей части.
\item Ответ --- последний принятый план.
\end{enumerate}

\begin{theorem}
Для любых параметров с $\pi(\mathbf5)\le\bar\varepsilon$ правило с порогом $\alpha P_1-\bar\varepsilon$ выбирает план с $\pi\le\alpha P_1$ с вероятностью не меньше $1-\gamma$ (с оценкой $\hat P_1$ вместо $P_1$ гарантия приближённая).
\end{theorem}

\begin{proof}
При фиксированной обучающей части путь неслучаен, а $S_k\sim\Bin(m,\rho(n^{(k)}))$, и $\rho$ монотонно убывает вдоль пути. Пусть $k^*$ --- наибольший номер с $\rho(n^{(k^*)})>\alpha P_1-\bar\varepsilon$. Ответ плохой, только если правило приняло $k^*$, то есть $U_\gamma(S_{k^*},m)\le\alpha P_1-\bar\varepsilon<\rho(n^{(k^*)})$; по свойству границы Клоппера--Пирсона вероятность этого не больше $\gamma$. Иначе $\pi=\rho+\pi(\mathbf5)\le\alpha P_1$.
\end{proof}

При $m=35$ и $\alpha\hat P_1-\bar\varepsilon\approx0.06$ правило принимает план, только если на проверочных релизах он не пропустил ни одного ($U_{0.1}(0,35)=0.064$, $U_{0.1}(1,35)=0.107$). Гарантия точная, но дорогая.

\section{Мир задачи}

\begin{center}\small
\begin{tabular}{@{}lcccc@{}}
\toprule
архетип & тестов & $s_i$ & $p_i$ & стоимость (множитель) \\
\midrule
идеальный & $6$ & $0.80$--$0.95$ & $0.022$--$0.050$ & $1$ \\
стеснительный & $6$ & $0.35$--$0.50$ & $0.022$--$0.050$ & $1$ \\
слабый & $3$ & $0.30$--$0.45$ & $0.022$--$0.050$ & $1$ \\
дорогой точный & $3$ & $0.90$--$0.99$ & $0.022$--$0.058$ & $10$ \\
дешёвый средний & $12$ & $0.35$--$0.55$ & $0.022$--$0.050$ & $0.3$ \\
не ломается & $5$ & $0.85$--$0.95$ & $0$--$0.0014$ & $1$ \\
часто ломается & $5$ & $0.45$--$0.65$ & $0.090$--$0.144$ & $1$ \\
\bottomrule
\end{tabular}
\end{center}

Базовая стоимость прогона --- логнормальная с медианой $8$ и $\sigma=0.4$, умноженная на множитель. Опубликованный мир --- seed $24$; его стоимости совпадают с уже выданным \texttt{evals.csv} (стоимости в генераторе зависят только от seed и долей архетипов). История --- $70$ релизов, валидация --- $30$. В мире $24$ доля релизов с проблемой $0.81$ (в истории $56$ из $70$, на валидации $25$ из $30$), в среднем $1.6$ проблемы на релиз; $\pi(\mathbf1)=0.207$ ($\pi_{\rm c}=0.256$), $\pi(\mathbf5)=0.012$ ($\pi_{\rm c}=0.015$); требование $\pi_{\rm c}\le0.1$ --- это $\pi\le0.081$.

\section{Что получается}

$60$ миров, по одной истории и одной валидации в каждом. Стоимость --- в долях от «все по $5$»; «пропущено» --- среднее число пропущенных релизов с проблемой на $30$ релизах валидации (из них в среднем $24$ проблемные; считаются видимые пропуски, $\Bin(30,\pi-\pi(\mathbf5))$, у плана с $\pi_{\rm c}=0.1$ ожидается около $2.1$). Столбец $\pi_{\rm c}$ --- истинная условная вероятность. Каждый метод целится в $0.1\cdot\hat P_1$ по своей оценке $\hat P_1$. Расчёт --- \texttt{evals\_task/numbers.py}.

\begin{center}\small
\begin{tabular}{@{}p{0.44\linewidth}cccc@{}}
\toprule
решение & $\pi_{\rm c}$ & \thead{миров с\\$\pi_{\rm c}\le0.1$} & стоимость & \thead{пропущено\\из $30$} \\
\midrule
все по разу & $0.248$ & $0$ & $0.20$ & $5.4$ \\
частоты, подстановка & $0.283$ & $0$ & $0.10$ & $6.1$ \\
общая $s$ для всех тестов & $0.186$ & $0$ & $0.12$ & $3.8$ \\
ОММ, подстановка & $0.138$ & $0.08$ & $0.17$ & $2.6$ \\
ОММ, дельта-метод & $0.114$ & $0.32$ & $0.21$ & $2.1$ \\
ОММ, бутстреп-квантиль без переотбора & $0.106$ & $0.42$ & $0.22$ & $1.9$ \\
ОММ, подстановка с целью $\alpha/2$ & $0.072$ & $0.88$ & $0.37$ & $1.1$ \\
ОММ, бутстреп всей процедуры, уровень $0.9$, $B=40$ & $0.074$ & $0.75$ & $0.38$ & $1.1$ \\
\textbf{ОММ, бутстреп всей процедуры, уровень $0.95$, $B=100$} & $0.067$ & $\mathbf{0.90}$ & $\mathbf{0.41}$ & $1.1$ \\
ОМП: подстановка / дельта / бутстреп & $0.04$--$0.05$ & $0.73$--$0.80$ & $0.78$ & $0.5$--$0.8$ \\
контроль $\E N$ / $\Prob(N\ge1)$ & $0.015$ & $1$ & $0.98$ / $0.97$ & $0$ \\
все по $2$ & $0.104$ & $0.25$ & $0.40$ & $2.3$ \\
отложенные релизы, Клоппер--Пирсон & $0.034$ & $0.98$ & $0.77$ & $0.4$ \\
оракул & $0.100$ & $1$ & $0.17$ & $1.9$ \\
все по $5$ & $0.015$ & $1$ & $1$ & $0$ \\
\bottomrule
\end{tabular}
\end{center}

\begin{itemize}
\item \textbf{Подстановки} не дороже оракула, но в требование почти не укладываются; частоты хуже нынешней практики --- тесты без событий не гоняются вовсе. Подстановка доли пропущенных релизов прямо по истории (без модели) переобучается: $\pi_{\rm c}=0.31$ при цели $0.1$.
\item \textbf{ОМП} вырождается и почти всё гоняет: надёжно, но дорого --- случайная перестраховка, а не метод.
\item \textbf{Контроль $\E N$ или $\Prob(N\ge1)$} вместо $\pi$ не экономит: подмена критерия стоит $0.97$--$0.98$ полной проверки. Та же подмена без модели --- доля релизов истории, где план упустил хоть одну проблему, --- даёт $\pi_{\rm c}=0.07$ при цене $0.37$: надёжно, но вдвое дороже оракула.
\item \textbf{Эталон} --- ОММ и бутстреп всей процедуры с уровнем $0.95$ и $B=100$: укладывается в $90\%$ миров, в $2.4$ раза дороже оракула и на $59\%$ дешевле полной проверки. С уровнем $0.9$ и $B=40$ бутстреп не дотягивает ($0.75$): при двух событиях на тест он недооценивает разброс. Грубая альтернатива --- наивная подстановка с целью $\alpha/2=0.05$: укладывается в $0.88$ миров при стоимости $0.37$; запас по бутстрепу в этом режиме не лучше запаса «вдвое», разница --- в обосновании.
\item \textbf{Отложенные релизы} дают точную гарантию за $0.77$.
\item \textbf{«Все по $2$»} при условном критерии уже не база: $\pi_{\rm c}=0.104$, укладывается лишь в четверти миров.
\item \textbf{Вырожденный план} «$15$ самых часто красных тестов по $3$ прогона» имеет $\pi_{\rm c}=0.28$ и проходит порог в $17\%$ выборок --- на опубликованной ему повезло.
\end{itemize}

\paragraph{Опубликованные данные.} Мир фиксирован стоимостями, поэтому выбиралась реализация истории и валидации: из $300$ реализаций --- та, на которой при безусловном критерии лестница решений была видна и по истине, и на валидации, а стоимость эталона близка к типичной для мира. При условном критерии цели всех методов сдвинулись, и на той же валидации наивная подстановка ОММ проходит порог; реализацию не меняли, потому что история уже выдана. Данные --- \texttt{data/B\_rel\_s24}; запасные реализации --- \texttt{data/B\_rel\_s24\_d146} и \texttt{data/B\_rel\_s24\_d153}; планы всех решений --- \texttt{hidden/plans}.

\begin{center}\small
\begin{tabular}{@{}p{0.40\linewidth}ccccc@{}}
\toprule
решение (мир $24$) & $\pi_{\rm c}$ & \thead{пропущено\\из $25$} & \thead{проходит\\порог $3$} & стоимость & прогонов \\
\midrule
все по разу & $0.256$ & $7$ & нет & $0.20$ & $40$ \\
частоты, подстановка & $0.277$ & $5$ & нет & $0.10$ & $52$ \\
общая $s$ & $0.213$ & $5$ & нет & $0.12$ & $56$ \\
ОМП, подстановка & $0.174$ & $4$ & нет & $0.15$ & $73$ \\
ОММ, подстановка & $0.166$ & $3$ & да & $0.15$ & $74$ \\
ОММ, дельта-метод & $0.122$ & $1$ & да & $0.20$ & $88$ \\
ОМП, дельта-метод & $0.126$ & $2$ & да & $0.19$ & $88$ \\
ОММ, бутстреп-квантиль & $0.122$ & $1$ & да & $0.20$ & $88$ \\
ОММ, подстановка с целью $\alpha/2$ & $0.080$ & $0$ & да & $0.30$ & $107$ \\
\textbf{ОММ, бутстреп всей процедуры (уровни $0.9$ и $0.95$ дали один план)} & $0.065$ & $0$ & да & $0.35$ & $115$ \\
все по $2$ & $0.108$ & $4$ & нет & $0.40$ & $80$ \\
контроль $\E N$ & $0.015$ & $0$ & да & $1$ & $200$ \\
отложенные релизы & $0.021$ & $0$ & да & $0.71$ & $191$ \\
оракул & $0.100$ & $0$ & да & $0.18$ & $86$ \\
\bottomrule
\end{tabular}
\end{center}

На $200$ историях этого мира эталон (уровень $0.95$) укладывается в требование в $87\%$ случаев ($82$--$91\%$) при средней стоимости $0.42$, бутстреп с уровнем $0.9$ --- в $76\%$ ($70$--$81\%$) при $0.36$; дельта-метод --- в $25\%$, наивная подстановка ОММ --- в $2.5\%$ (средняя цена $0.18$).

\section{Как оценивать ответы}

Ответ команды --- \texttt{plan.csv} (\texttt{eval\_id}, \texttt{n\_runs}). На $30$ скрытых релизах (\texttt{dataset\_valid.csv}) записаны все пять прогонов; релиз с проблемой --- тот, где хоть у одного теста среди пяти записей есть красный. Считаются (\texttt{evals\_task/score.py}):
\begin{itemize}
\item \textbf{пропущено} --- релизы с проблемой, на которых первые $n_i$ записей всех тестов зелёные;
\item \textbf{потрачено} --- $\sum_ic_in_i$.
\end{itemize}
План проходит, если пропустил не больше $3$ релизов; прошедшие ранжируются по стоимости, не прошедшие --- после них. Считаются видимые пропуски, поэтому их число --- $\Bin(30,\rho)$ с $\rho=\pi_{\rm c}P_1-\pi(\mathbf5)=0.81\,\pi_{\rm c}-0.012$; у плана с $\pi_{\rm c}=0.1$ среднее $2.1$, и порог $3$ --- его $80\%$-квантиль. Вероятность пройти порог $3$: при $\pi_{\rm c}=0.05$ --- $0.99$, $0.08$ --- $0.93$, $0.10$ --- $0.85$, $0.13$ --- $0.69$, $0.15$ --- $0.58$, $0.20$ --- $0.32$, $0.30$ --- $0.06$. То есть порог надёжно отсекает лишь $\pi_{\rm c}\gtrsim0.2$; между $0.1$ и $0.17$ исход --- в значительной мере удача.

\paragraph{Разбор.} Решено оставить $30$ релизов и раскрыть истину на разборе: для каждого плана команды показать истинную условную вероятность пропуска (\texttt{python3 -m evals\_task.score ... --exact}, столбец P(miss|prob)) рядом с числом пропущенных и честно сказать, что на $30$ релизах между $\pi_{\rm c}=0.1$ и $0.17$ исход --- во многом удача. Примеры на опубликованных данных: наивная подстановка ОММ с целью $0.1$ даёт $\pi_{\rm c}=0.166$, но пропускает лишь $3$ из $25$ и проходит порог дешевле оракула --- такая удача у неё случается примерно в половине выборок; план «$15$ самых часто красных тестов по $3$ прогона» пропускает $2$ из $25$ и стоит $0.13$, а по истине пропускает $29\%$ проблемных релизов; подстановка доли пропущенных релизов по истории с целью $0.1$ переобучается --- $8$ пропусков, $\pi_{\rm c}=0.33$; подстановка $\Prob(N\ge1)$ по истории даёт $\pi_{\rm c}=0.131$ и $5$ пропусков.

\paragraph{Что валидация различает.} Дешёвые и грубые решения (частоты, общая $s$, «все по разу») пропускают $4$--$8$ релизов и отделяются уверенно. Наивная подстановка ОММ превышает $\alpha$ в $1.3$--$1.7$ раза --- на $30$ релизах это видно лишь в части миров, и на опубликованной реализации при условном критерии не видно; честное сравнение методов --- по истинным параметрам на многих мирах (таблица выше).

\end{document}
